\documentclass[12pt,a4paper]{article}
\usepackage[UTF8]{ctex}        % 中文支持
\usepackage{geometry}
\geometry{left=2.54cm,right=2.54cm,top=3.18cm,bottom=3.18cm}
\usepackage{enumerate}        % 列表格式
\usepackage{titlesec}         % 标题格式
\usepackage{xcolor}           % 颜色包
\usepackage{indentfirst}
\setlength{\parindent}{2em}
\usepackage{enumitem}
\usepackage{enumerate}


\usepackage[
    colorlinks=true,
    linkcolor=red,    % 目录、图表、章节引用：黑色
    citecolor=red,      % \cite{} 参考文献引用：红色（你需要的）
    urlcolor=red     % 网址链接：黑色
]{hyperref}

\usepackage{listings}
\lstset{
    language=Matlab,          % 指定语言为Matlab
    basicstyle=\ttfamily\small,% 代码基础字体
    keywordstyle=\color{blue}, % function/end 关键字蓝色
    commentstyle=\color{green!60!black}, % %注释绿色
    stringstyle=\color{red},   % 字符串红色
    numbers=left,              % 左侧行号（不需要写 numbers=none）
    numberstyle=\tiny\color{gray},
    frame=single,              % 外框单线
    backgroundcolor=\color{gray!8}, % 浅灰背景
    tabsize=4,                 % Tab缩进4空格
    showstringspaces=false,    % 不显示空格下划线
    breaklines=true,           % 长代码自动换行
    escapeinside={/*@}{@*/},   % 代码内可插LaTeX命令
}

% 字体设置

% ===================== 自定义功能区 =====================
\newcommand{\textred}[1]{\textcolor{red}{#1}}           %红色字体
\newcommand{\textgreen}[1]{\textcolor{green}{#1}}       %绿色字体
\newcommand{\textblue}[1]{\textcolor{blue}{#1}}         %绿色字体
\newcommand{\textyellow}[1]{\textcolor{yellow}{#1}}     %黄色字体
% \newcommand{\wordnote}[2]{#1\,\raisebox{0.25ex}{\scriptsize #2}}
\newcommand{\upmean}[1]{\raisebox{0.25ex}{\scriptsize #1}} %标注





%========================== 正文 ============================


\begin{document}

\section{计算机视觉中的图像特征}
\subsection{填写函数的方式和原因}

\textbf{定义函数}

函数定义由三个部分组成。

1. 函数声明以关键字 function 开头，定义调用函数的语法。

2. 函数体包含使用输入计算函数输出的命令。输出必须在函数体中定义。

3. end 关键字是函数的最后一行。

\begin{lstlisting}
function out = myFunction(input)
% code that calculates out
end
\end{lstlisting}


\textbf{调用优先级}

在 MATLAB 中，设置了规则以解释任何命名项。这些规则称为函数优先顺序。
可以使用以下顺序解决最常见的引用冲突：

变量

当前脚本中定义的函数

当前文件夹中的文件

MATLAB 搜索路径上的文件


\subsection{计算机视觉入门之旅}

\href{https://matlabacademy.mathworks.com/cn/details/computer-vision-onramp/orcv}{计算机视觉入门之旅}

课程简介：通过将典型工作流（检测跟踪）应用于海龟爬向大海的视频，
学习计算机视觉的基础知识。您将了解特征在计算机视觉中的作用、如
何标记数据、训练目标检测器以及在视频中跟踪野生动物。

{\large 使用目标检测跟踪海龟}

在本课程中，您将学习如何使用目标检测来跟踪视频中的目标。课程将从导入视频数据开始，
逐步讲解此检测跟踪工作流。

Import Video Data $\Rightarrow $ 
Label Ground Truth Data $\Rightarrow $ 
Train Detector$\Rightarrow $ 
Detect Objects in Each Frame$\Rightarrow $ 
Track Objects

\subsubsection{创建视频读取器}

可以使用 VideoReader 函数创建一个视频读取器，用于从文件中读取视频数据。

\begin{lstlisting}
v = VideoReader("myVid.mp4")
\end{lstlisting}

创建视频读取器后，您可以使用 readFrame 函数从视频文件中提取帧。

\begin{lstlisting}
v = VideoReader(filename)
fr = readFrame(v);
\end{lstlisting}

readFrame 的输出是图像，它是由像素值组成的数组。您可以在 readFrame 
函数结尾使用分号，以避免 MATLAB 显示数百万个数字。

视频的一个帧就是一个图像。从视频文件中提取帧后，您可以像处理任何图像
一样处理该帧。

\begin{lstlisting}
imshow(frame)
%显示视频某一帧
\end{lstlisting}

在输出窗格的顶部，您可以看到 turtleVideo 的一些属性。例如，FrameRate 
属性指示此视频的帧速率，即每秒显示的帧数

您可以用圆点表示法提取这些属性。

\begin{lstlisting}
prop = obj.PropertyName
\end{lstlisting}

在上一个任务中，t 是 0.0333，即 $\frac{1}{frame~rate}$
 秒。在第 1 行上创建视频读取器时，CurrentTime 属性的值为 0。

您可以使用 imshowpair 函数叠加两个图像，并以绿色和品红色显示它们的差异。

\begin{lstlisting}
imshowpair(im1,im2)
\end{lstlisting}

读取特定帧

您可以反复调用 readFrame 以读取此视频的任一帧。但如何直接跳到第 90 帧呢？
或如何回退到某个更早的帧？在本练习中，您将学习如何读取特定帧以及如何导航
到视频中的任一帧。

在输出窗格的顶部，您可以在 turtleVideo 的输出中看到视频的总帧数和持续时间。
变量 tMid 大约在视频的一半位置。

\begin{lstlisting}
turtleVideo = 
  VideoReader - 属性:

   常规属性:
            Name: 'turtles.avi'
            Path: '/users/mwa0000035781096/.training/.work'
        Duration: 3.3333
     CurrentTime: 0
       NumFrames: 100

   视频属性:
           Width: 400
          Height: 600
       FrameRate: 30
    BitsPerPixel: 24
     VideoFormat: 'RGB24'
\end{lstlisting}

\subsection{从参考图像中检测并提取特征}

识别海龟的特征

要进行目标检测，您需要具有待检测目标的参考图像。参考图像应仅包含对检测
该目标有用的特征。在此任务中，您将通过在视频帧中裁剪海龟周围的区域，来
创建一个海龟的参考图像。

要指定要裁剪的矩形，需提供左上角的坐标以及矩形的宽度和高度（以像素为单
位）。对于最小的 x 值和 y 值，您可以分别使用图像的列数和行数。

您可以通过 imcrop 函数指定要保留的矩形来裁剪图像。
imCropped = imcrop(im,rect);
将矩形指定为一个四元素向量 [xmin ymin width height]。

\begin{lstlisting}
    rect = [xmin ymin width height];
    imCropped = imcrop(im,rect);
\end{lstlisting}

您可以从矩阵或灰度图像中提取特征。由于 turtle 和 frame 是彩色图像，提取特征之前需将其转换为灰度图像。

您可以使用 im2gray 函数基于彩色图像创建灰度图像。

\begin{lstlisting}
imGS = im2gray(im);
\end{lstlisting}

识别特征过程分为两部分：检测和提取。要检测图像中的类连通域特征，您可以使用 
detectSIFTFeatures 函数。

输出是一系列感兴趣点的列表，其中可能包含尺度不变特征变换 (SIFT) 特征。

\begin{lstlisting}
pts = detectSIFTFeatures(im)
\end{lstlisting}

现在您有了感兴趣点，可以使用它们来提取特征。感兴趣点标识了可查找特征的位置；
一个点附近可能有多个特征，也可能一个特征也没有。

您可以对任何类型的特征使用 extractFeatures 函数。

第一个输出是特征矩阵。第二个输出是另一个 SIFTpoints 数组，其中包含这些特征的具体位置。

\begin{lstlisting}
[feat,ftPts] = extractFeatures(im,pts)
\end{lstlisting}

您可以通过绘图对提取的特征可视化。大多数特征对应于海龟图像的某个区域。
SIFT 特征是一种连通域特征。其他特征类型对应使用相同语法的函数。
使用的特征类型取决于图像中的特征以及使用方法。

\begin{lstlisting}
imshow(turtle)
hold on
plot(featurePoints)
hold off
\end{lstlisting}

您已从参考图像和原始视频帧中提取了特征，现在可以使用 matchFeatures 
函数来匹配这些特征。

输出是一个包含匹配特征的索引的两列矩阵。indexedPairs 的第一列包含 
feat1 的索引。第二列包含 feat2 的索引

\begin{lstlisting}
indexedPairs = matchFeatures(feat1,feat2)
\end{lstlisting}



\end{document}